\otazka{22. Ochrana osobních údajů}{Právní úprava. Principy a zásady zpracování osobních údajů - definice, základní zásady, posouzení rizik, test proporcionality. Účely zpracování. Právní tituly. Zákon o zpracování osobních údajů - ÚOOÚ a jeho úloha a postavení v České republice.}{BI301K}

\subsection*{22. Ochrana osobních údajů -- stručněji}
GDPR není zákaz práce s daty, ale pravidlový rámec. Správná odpověď nezačíná souhlasem, ale účelem: proč údaje sbírám, jaký titul mi to dovoluje, jaké zásady dodržím a jak doložím odpovědnost.

\textbf{1. Právní rámec a pojmy}
\begin{itemize}
  \item Základ je GDPR jako přímo použitelné nařízení EU a český zákon o zpracování osobních údajů. Zvláštní předpisy doplňují pravidla pro zdravotnictví, práci, elektronické komunikace nebo policii.
  \item Osobní údaj je informace o identifikované nebo identifikovatelné fyzické osobě. Správce určuje účely a prostředky, zpracovatel jedná pro správce podle jeho pokynů.
  \item Pseudonymizace není anonymizace. Pseudonymizovaná data zůstávají osobními údaji, pokud existuje vazba zpět na člověka.
\end{itemize}

\textbf{2. Zásady zpracování}
\begin{itemize}
  \item Zásady jsou zákonnost, korektnost, transparentnost, účelové omezení, minimalizace, přesnost, omezení uložení, integrita a důvěrnost a odpovědnost.
  \item Odpovědnost znamená, že správce musí soulad nejen dodržet, ale také doložit: záznamy, smlouvy, DPIA, bezpečnostní opatření, interní pravidla a školení.
\end{itemize}

\textbf{3. Účel a právní titul}
\begin{itemize}
  \item Účel je skutečně to, \textbf{proč} údaje sbíráte. Musí být jasný, konkrétní a stanovený předem. Právní titul je až zákonný \textbf{důvod}, který dovoluje tento účel uskutečnit.
  \item Tituly podle čl. 6 jsou souhlas, smlouva, právní povinnost, životně důležitý zájem, veřejný zájem/výkon veřejné moci a oprávněný zájem.
  \item Souhlas není univerzální řešení. Musí být svobodný, konkrétní, informovaný, jednoznačný a odvolatelný.
\end{itemize}

\textbf{4. Zvláštní kategorie, oprávněný zájem a DPIA}
\begin{itemize}
  \item Zvláštní kategorie údajů jsou citlivé údaje, například zdravotní, biometrické pro identifikaci, genetické, politické, náboženské nebo údaje o sexuální orientaci. Potřebují titul podle čl. 6 a výjimku podle čl. 9.
  \item Oprávněný zájem vyžaduje balancing test: oprávněný cíl, nezbytnost a vyvážení vůči právům subjektu údajů.
  \item DPIA se dělá před rizikovým zpracováním, například rozsáhlé profilování, biometrika, rozsáhlé kamerové sledování nebo zpracování citlivých údajů.
\end{itemize}

\textbf{5. Práva subjektů, DPO a ÚOOÚ}
\begin{itemize}
  \item Subjekty mají právo na informace, přístup, opravu, výmaz, omezení, přenositelnost, námitku a ochranu proti čistě automatizovanému rozhodování.
  \item DPO radí, kontroluje soulad, školí a komunikuje s ÚOOÚ, ale odpovědnost za GDPR zůstává správci.
  \item ÚOOÚ je český dozorový úřad: kontroluje, řeší stížnosti, ukládá opatření a pokuty a spolupracuje v evropském systému.
\end{itemize}

\subsection*{Kostra odpovědi}
\begin{enumerate}
  \item Říct právní rámec: GDPR, zákon o zpracování osobních údajů, zvláštní předpisy.
  \item Vymezit základní pojmy: osobní údaj, subjekt údajů, správce, zpracovatel, zpracování.
  \item Vyjmenovat zásady čl. 5 GDPR.
  \item Vysvětlit účel a právní titul: účel je proč, titul je proč to smím.
  \item Vyjmenovat právní tituly čl. 6 a zvláštní kategorie čl. 9.
  \item Popsat riziko, DPIA, test proporcionality a DPO.
  \item Zmínit práva subjektů, porušení zabezpečení, ÚOOÚ a předávání mimo EHP.
\end{enumerate}

\subsection*{Právní úprava}
\begin{itemize}
  \item \textbf{GDPR} - hlavní přímo použitelné nařízení EU pro ochranu osobních údajů.
  \item \textbf{Zákon o zpracování osobních údajů} - český adaptační zákon; doplňuje GDPR, řeší ÚOOÚ, některé výjimky a zvláštní režimy.
  \item \textbf{Zvláštní předpisy} - zdravotnictví, pracovní právo, školství, banky, elektronické komunikace, obchodní sdělení, policie.
  \item \textbf{EDPB} - Evropský sbor pro ochranu osobních údajů sjednocuje výklad v EU.
\end{itemize}

\subsection*{Základní pojmy}
\begin{itemize}
  \item \textbf{Osobní údaj} - informace o identifikované nebo identifikovatelné fyzické osobě: jméno, e-mail, rodné číslo, IP adresa, cookie, poloha, fotka, hlas, kombinace údajů.
  \item \textbf{Subjekt údajů} - člověk, kterého se údaje týkají.
  \item \textbf{Správce} - určuje, proč a jak se údaje zpracovávají.
  \item \textbf{Zpracovatel} - zpracovává údaje pro správce podle jeho pokynů, například cloud, účetní, mailingový nástroj.
  \item \textbf{Zpracování} - sběr, uložení, použití, předání, změna, zveřejnění, výmaz.
  \item \textbf{Pseudonymizace} - vazba na osobu stále existuje, GDPR platí.
  \item \textbf{Anonymizace} - osobu už nelze rozumně znovu identifikovat, GDPR se neuplatní.
\end{itemize}

\subsection*{Zásady zpracování}
\begin{longtable}{p{0.30\linewidth}p{0.62\linewidth}}
\toprule
\textbf{Zásada} & \textbf{Co říct}\\
\midrule
Zákonnost, korektnost, transparentnost & Musím mít právní titul a člověk má vědět, co s údaji dělám.\\
Účelové omezení & Údaje sbírám pro jasný účel a nepoužiji je libovolně jinak.\\
Minimalizace & Sbírám jen údaje, které opravdu potřebuji.\\
Přesnost & Údaje mají být správné a aktualizované.\\
Omezení uložení & Nedržím údaje déle, než je potřeba.\\
Integrita a důvěrnost & Musím zajistit bezpečnost: přístupy, šifrování, logy, zálohy.\\
Odpovědnost & Správce musí soulad nejen dodržet, ale i doložit.\\
\bottomrule
\end{longtable}

\subsection*{Účely a právní tituly}
\pozor{Účel je skutečně to, \textbf{proč} údaje sbírám. Musí být jasný, konkrétní a stanovený předem. Právní titul je až zákonný \textbf{důvod}, který dovoluje tento účel uskutečnit.}

\begin{longtable}{p{0.28\linewidth}p{0.62\linewidth}}
\toprule
\textbf{Právní titul čl. 6} & \textbf{Příklad}\\
\midrule
Souhlas & Newsletter pro nového zájemce; musí být svobodný, konkrétní, informovaný a odvolatelný.\\
Plnění smlouvy & Jméno a adresa pro doručení objednávky.\\
Právní povinnost & Uchování faktur podle účetních a daňových pravidel.\\
Životně důležitý zájem & Zpracování údajů pacienta v akutní nouzi.\\
Veřejný zájem / veřejná moc & Matrika, úřad, výkon veřejné správy.\\
Oprávněný zájem & Bezpečnostní logy, vymáhání pohledávky, ochrana majetku, některý přímý marketing.\\
\bottomrule
\end{longtable}

\subsection*{Oprávněný zájem}
Oprávněný zájem není kouzelné zaklínadlo. Správce má udělat \textbf{LIA / balancing test}:
\begin{enumerate}
  \item \textbf{Účel} - mám konkrétní a legitimní zájem?
  \item \textbf{Nezbytnost} - nejde toho dosáhnout méně invazivně?
  \item \textbf{Vyvážení} - nepřevažují práva a rozumná očekávání člověka?
\end{enumerate}

\priklad{Kamera u vstupu do skladu kvůli krádežím může obstát: jasný účel, krátká retence, záběr jen na vstup, informování osob. Kamera se zvukem nad celým pracovištěm kvůli kontrole výkonu zaměstnanců už pravděpodobně neobstojí.}

\subsection*{Zvláštní kategorie údajů}
Zvláštní kategorie podle čl. 9 GDPR jsou citlivé údaje: zdraví, genetické údaje, biometrika pro jedinečnou identifikaci, sexuální orientace, politické názory, náboženství, členství v odborech, rasový nebo etnický původ.

U těchto údajů nestačí běžný titul. Potřebuji:
\begin{itemize}
  \item \textbf{čl. 6} - obecný právní titul,
  \item \textbf{čl. 9} - výjimku ze zákazu zpracování.
\end{itemize}

Typické výjimky podle čl. 9: výslovný souhlas, pracovní právo a sociální zabezpečení, životně důležité zájmy, zdravotní péče, veřejné zdraví, významný veřejný zájem, právní nároky, výzkum a archivace. Údaje o trestních odsouzeních a trestných činech mají zvláštní režim podle čl. 10.

\subsection*{Riziko, DPIA a proporcionalita}
GDPR je postavené na riziku. Čím vyšší riziko pro člověka, tím silnější opatření. \textbf{DPIA} se dělá před zahájením zpracování, pokud zpracování pravděpodobně povede k vysokému riziku.

DPIA typicky řeším u:
\begin{itemize}
  \item rozsáhlého profilování s významnými účinky,
  \item rozsáhlého zpracování zvláštních kategorií údajů,
  \item rozsáhlého kamerového sledování veřejně přístupných prostor,
  \item biometriky, AI, nových technologií nebo spojování velkých datových sad.
\end{itemize}

Test proporcionality řekni čtyřmi kroky: legitimní cíl, vhodnost, nezbytnost, přiměřenost v užším smyslu.

\subsection*{DPO, práva a incidenty}
\textbf{DPO / pověřenec} radí, kontroluje, školí, komunikuje s ÚOOÚ a je kontaktním místem. Není to osoba, která přebírá odpovědnost správce. Povinný je hlavně u veřejné moci, rozsáhlého systematického monitorování a rozsáhlého zpracování citlivých údajů.

Subjekt údajů má práva: informace, přístup, oprava, výmaz, omezení, přenositelnost, námitka, nebýt předmětem čistě automatizovaného rozhodování s významnými účinky.

\textbf{Námitka} je důležitá hlavně u oprávněného zájmu a veřejného zájmu. U přímého marketingu platí: když člověk vznese námitku, marketing musí skončit. \textbf{Výmaz} není absolutní: faktury nebo údaje nutné pro právní nároky se nemažou jen proto, že si to někdo přeje.

Porušení zabezpečení se hlásí ÚOOÚ zpravidla do 72 hodin, pokud hrozí riziko pro práva a svobody. Subjektům údajů se oznamuje při vysokém riziku.

\subsection*{ÚOOÚ a předávání mimo EHP}
ÚOOÚ je nezávislý dozorový úřad: řeší stížnosti, kontroly, opatření, pokuty, metodiku a spolupráci v EDPB. Sankce mohou u nejzávažnějších porušení dosáhnout až 20 mil. eur nebo 4 \% celosvětového obratu.

Předávání mimo EHP je možné hlavně přes rozhodnutí o odpovídající ochraně, standardní smluvní doložky, závazná podniková pravidla nebo výjimky. U standardních doložek je potřeba myslet i na rizika práva cílové země.

\zkouska{Souvisle: u GDPR nejdřív řeknu právní rámec, potom pojmy, zásady, účel a právní titul. Nikdy nezačínám automaticky souhlasem. U citlivých údajů potřebuji čl. 6 i čl. 9. U oprávněného zájmu udělám balancing test. U rizikového zpracování udělám DPIA před spuštěním. Správce musí umět soulad doložit a ÚOOÚ nad tím vykonává dohled.}

\zdroje{GDPR 2016/679; zákon o zpracování osobních údajů; zákon o některých službách informační společnosti; metodiky ÚOOÚ a EDPB; přiložené PDF a starší výpisky.}

\newpage

\iffalse
\otazka{22. Ochrana osobních údajů}{Právní úprava. Principy a zásady zpracování osobních údajů - definice, základní zásady, posouzení rizik, test proporcionality. Účely zpracování. Právní tituly. Zákon o zpracování osobních údajů - ÚOOÚ a jeho úloha a postavení v České republice.}{BI301K}

\subsection*{Kostra odpovědi}
\begin{enumerate}
  \item Právní rámec: GDPR jako přímo použitelné nařízení, zákon o zpracování osobních údajů, zvláštní předpisy a ePrivacy/obchodní sdělení.
  \item Pojmy: osobní údaj a identifikovatelnost, subjekt údajů, správce, zpracovatel, příjemce, pseudonymizace vs anonymizace.
  \item Zvláštní kategorie údajů: čl. 6 a čl. 9 zároveň, čl. 10 pro trestní údaje.
  \item Zásady čl. 5 na příkladech: zákonnost, transparentnost, účel, minimalizace, přesnost, omezení uložení, integrita/důvěrnost, odpovědnost.
  \item Nejdřív účel, potom právní titul: šest titulů čl. 6, souhlas, soft opt-in a oprávněný zájem.
  \item Oprávněný zájem: LIA/balancing test - oprávněný cíl, nezbytnost, vyvážení zájmů a právo námitky.
  \item DPIA, proporcionalita, privacy by design/default a DPO.
  \item Práva subjektů: přístup, oprava, výmaz, omezení, přenositelnost, námitka, automatizované rozhodování.
  \item Porušení zabezpečení, ÚOOÚ, EDPB, one-stop-shop a předávání mimo EHP.
\end{enumerate}

\subsection*{Právní úprava}
Základem je GDPR, tedy nařízení EU 2016/679, použitelné od 25. května 2018. Je přímo použitelné a sjednocuje pravidla ochrany osobních údajů v EU. To je rozdíl proti směrnicím typu NIS2: GDPR se netransponuje, ale přímo platí. Český zákon o zpracování osobních údajů (č. 110/2019 Sb.) je adaptační zákon. GDPR nenahrazuje, jen doplňuje tam, kde nařízení nechává prostor členskému státu.

Český zákon řeší zejména postavení ÚOOÚ, některé výjimky pro žurnalistiku, vědu, archivnictví a svobodu projevu, samostatný režim pro zpracování v trestněprávní oblasti a také věkovou hranici dítěte pro souhlas u služeb informační společnosti: v ČR 15 let. Vedle toho existují zvláštní předpisy pro zdravotnictví, pracovní právo, elektronické komunikace, obchodní sdělení, bankovnictví, školství, archivnictví, policii a zpravodajské služby.

GDPR není "zákaz zpracování dat". Je to pravidlový rámec: osobní údaje zpracovávat lze, ale musí být jasný účel, právní titul, přiměřenost, bezpečnost a odpovědnost.

\subsubsection*{GDPR, trestněprávní oblast a územní působnost}
GDPR se neuplatní na všechno zpracování osobních údajů veřejnou mocí. Pokud příslušné orgány zpracovávají osobní údaje pro účely prevence, vyšetřování, odhalování nebo stíhání trestných činů, výkonu trestů nebo ochrany před hrozbami pro veřejnou bezpečnost, nastupuje zvláštní režim směrnice 2016/680, tzv. Law Enforcement Directive, a český zákon o zpracování osobních údajů. Stejný orgán proto může být jednou v režimu GDPR a jindy v režimu LED podle účelu. Policie při personalistice nebo nákupu služeb řeší GDPR; při vyšetřování trestného činu řeší zvláštní trestněprávní režim.

Územní působnost GDPR není omezena jen na firmy fyzicky v EU. GDPR se použije na zpracování v souvislosti s provozovnou správce nebo zpracovatele v EU. Vedle toho dopadá i na subjekty mimo EU, pokud nabízejí zboží nebo služby osobám v EU nebo monitorují jejich chování v EU, například cíleným trackingem, profilováním nebo behaviorální reklamou. Neznamená to, že každý web na světě automaticky spadá pod GDPR; rozhoduje reálné cílení na osoby v EU nebo monitorování jejich chování.

\pozor{Domácí výjimka se týká čistě osobní nebo domácí činnosti, například soukromého adresáře nebo rodinné fotografie. Neplatí automaticky, pokud osoba zveřejňuje údaje neurčité veřejnosti na internetu, provozuje veřejnou stránku, systematicky monitoruje okolí nebo zasahuje do práv cizích osob.}

\subsection*{Základní pojmy}
\textbf{Osobní údaj} je informace o identifikované nebo identifikovatelné fyzické osobě. Identifikovatelnost znamená, že osobu lze určit přímo nebo nepřímo pomocí rozumně dostupných prostředků. Jméno a rodné číslo jsou zřejmé příklady, ale osobním údajem může být i IP adresa, cookie identifikátor, poloha, hlas, fotografie, pracovní identifikátor nebo kombinace údajů.

Ne každý technický identifikátor je osobním údajem vždy a pro každého. Dynamická IP adresa nebo cookie je osobní údaj tehdy, pokud existují rozumně pravděpodobné prostředky, jak ji spojit s osobou. Pro poskytovatele služby to často osobní údaj bude; pro subjekt bez možnosti vazby to tak být nemusí.

\textbf{Subjekt údajů} je fyzická osoba, které se údaje týkají. \textbf{Správce} určuje účely a prostředky zpracování, tedy proč a jak se údaje zpracovávají. \textbf{Zpracovatel} zpracovává údaje pro správce a podle jeho pokynů, typicky cloud, externí účetní, mailingový nástroj nebo IT dodavatel; vztah se upravuje zpracovatelskou smlouvou podle čl. 28 GDPR. \textbf{Příjemce} je ten, komu jsou údaje zpřístupněny. \textbf{Zpracování} je skoro jakákoli operace s údaji: sběr, ukládání, změna, použití, předání, nahlížení, zveřejnění nebo výmaz.

\textbf{Pseudonymizace} nahrazuje identifikátory kódem, ale klíč nebo vazba existuje odděleně, takže GDPR stále platí. \textbf{Anonymizace} musí být prakticky nevratná; osobu už nelze rozumně pravděpodobnými prostředky znovu identifikovat, vyčlenit ani propojit. Pouhé odstranění jména není anonymizace, pokud zbylá data pořád umožní re-identifikaci.

\pozor{Pseudo není anonym. Pseudonymizovaná data zůstávají osobními údaji, protože existuje vazba zpět na člověka.}

\subsection*{Zvláštní kategorie údajů}
Zvláštní kategorie podle čl. 9 GDPR jsou citlivější osobní údaje: zdravotní údaje, genetické údaje, biometrické údaje zpracovávané za účelem jedinečné identifikace, údaje o sexuální orientaci, politických názorech, náboženském nebo filozofickém přesvědčení, členství v odborech a rasovém nebo etnickém původu. Jejich zpracování je v zásadě zakázané a teprve výjimka podle čl. 9 zákaz prolomí.

U zvláštních kategorií je potřeba splnit dvě vrstvy najednou:
\begin{itemize}
  \item \textbf{čl. 6} - obecný právní titul, tedy proč osobní údaj vůbec smím zpracovat,
  \item \textbf{čl. 9} - zvláštní výjimka, která dovolí zpracovat citlivý údaj.
\end{itemize}
Splnit jen jednu vrstvu nestačí. Nemocnice například zpracovává zdravotní údaje na základě právní povinnosti nebo poskytování zdravotní péče podle čl. 6 a současně používá výjimku podle čl. 9 pro zdravotní péči. Údaje o odsouzeních a trestných činech mají vlastní režim v čl. 10.

\priklad{Fitness aplikace sbírá počet kroků a tepovou frekvenci. Pokud z toho vyvozuje zdravotní stav nebo používá data pro zdravotní doporučení, může se dostat ke zvláštní kategorii údajů. Pak nestačí "uživatel souhlasil s podmínkami"; je potřeba přesně vyřešit titul podle čl. 6 a výjimku podle čl. 9, často výslovný souhlas.}

\priklad{Aktuální příklad z praxe - biometrická docházka: ÚOOÚ v kontrolách ukončených v roce 2025 řešil zaměstnavatele, kteří používali docházkové systémy založené na otiscích prstů nebo jejich šablonách. Potvrzená pointa je, že biometrický údaj je zvláštní kategorií podle GDPR, pokud se zpracovává za účelem jedinečné identifikace osoby. U zaměstnanců je navíc problematická svobodnost souhlasu, protože vztah je nerovný. Ke zkoušce: biometrika v práci není jen "moderní docházka"; vyžaduje čl. 6, výjimku podle čl. 9, minimalizaci, posouzení alternativ a často DPIA.}

\subsection*{Zásady zpracování}
GDPR stojí na zásadách:
\begin{itemize}
  \item \textbf{zákonnost, korektnost a transparentnost} - zpracování musí mít právní titul a člověk má srozumitelně vědět, co se děje; e-shop musí jasně informovat, že e-mail použije pro objednávku,
  \item \textbf{účelové omezení} - data se sbírají pro konkrétní účel a nelze je libovolně použít jinak; e-mail z objednávky nejde bez dalšího použít k profilování reklam,
  \item \textbf{minimalizace} - sbírat jen to, co je nezbytné; pro newsletter stačí e-mail, ne rodné číslo,
  \item \textbf{přesnost} - údaje mají být aktuální a správné; špatnou doručovací adresu je třeba opravit,
  \item \textbf{omezení uložení} - neuchovávat déle, než je potřeba; faktury podle účetních lhůt, marketingový kontakt jen do odvolání souhlasu nebo námitky,
  \item \textbf{integrita a důvěrnost} - zajistit bezpečnost: řízení přístupů, šifrování, zálohy, logování,
  \item \textbf{odpovědnost} - správce musí soulad nejen dodržet, ale i doložit: záznamy o činnostech, testy, DPIA, smlouvy, interní pravidla.
\end{itemize}

\subsection*{Účely a právní tituly}
Účel říká, proč údaje zpracovávám. Právní titul říká, proč to smím. Správné pořadí je: nejdřív účel, potom právní titul, potom rozsah údajů, doba uložení, informace subjektu a zabezpečení. Titul se nemá vybírat dodatečně jako náplast, když se původní důvod nehodí.

Čl. 6 GDPR zná šest titulů:
\begin{enumerate}
  \item \textbf{souhlas} - newsletter pro nového zájemce, pokud není jiný titul,
  \item \textbf{plnění smlouvy} - jméno a adresa pro doručení objednávky,
  \item \textbf{právní povinnost} - uchování účetních a daňových dokladů,
  \item \textbf{ochrana životně důležitých zájmů} - péče o pacienta v bezvědomí,
  \item \textbf{úkol ve veřejném zájmu nebo výkon veřejné moci} - matrika, veřejná správa,
  \item \textbf{oprávněný zájem} - bezpečnostní monitoring, ochrana majetku, vymáhání pohledávky, některý přímý marketing.
\end{enumerate}

Souhlas není univerzální náplast. Musí být svobodný, konkrétní, informovaný, jednoznačný a odvolatelný stejně snadno, jako byl udělen. Problém vzniká, když je služba podmíněna souhlasem se zpracováním, které pro službu není nezbytné: například nákup v e-shopu podmíněný souhlasem s marketingem.

U přímého elektronického marketingu je vedle souhlasu důležitý i český režim obchodních sdělení podle zákona o některých službách informační společnosti (č. 480/2004 Sb.) a tzv. soft opt-in: stávajícímu zákazníkovi lze za splnění podmínek nabízet vlastní obdobné výrobky nebo služby, pokud měl možnost jednoduše a zdarma odmítnout při získání kontaktu i v každé další zprávě.

\subsection*{Oprávněný zájem a balancing test}
Oprávněný zájem podle čl. 6 odst. 1 písm. f GDPR není volná karta. Správce musí před zpracováním projít a zdokumentovat test oprávněného zájmu, často označovaný jako LIA. Má tři kroky:
\begin{enumerate}
  \item \textbf{Test účelu / oprávněnosti cíle} - je zájem konkrétní, skutečný a zákonný? Typicky ochrana majetku, bezpečnost sítě, prevence podvodů, vymáhání pohledávky, přímý marketing vůči vlastním zákazníkům.
  \item \textbf{Test nezbytnosti} - je zpracování pro daný cíl skutečně potřebné a nejde stejného výsledku dosáhnout méně invazivně? Kamera na vstup může být nezbytná, kamera na celou jídelnu nebo šatnu už ne.
  \item \textbf{Balancing test / vyvážení} - nepřevažují práva, svobody a rozumná očekávání subjektu nad zájmem správce? Hodnotí se rozsah údajů, citlivost, vztah se subjektem, možnost námitek, dopad, transparentnost, záruky a bezpečnostní opatření.
\end{enumerate}

Výsledek má být písemně doložen, protože správce nese odpovědnost. Test oprávněného zájmu není totéž co DPIA a není totéž co ústavní test proporcionality. LIA řeší, zda mohu použít konkrétní právní titul. DPIA řeší rizika rizikového zpracování jako celek. Ústavní test proporcionality je obecnější test zásahu do základních práv, typicky u veřejné moci.

\pravobox{TEST OPRÁVNĚNÉHO ZÁJMU A JEHO HRANICE: Firma chce nasadit kameru u vstupu do skladu kvůli krádežím. Cíl je oprávněný: ochrana majetku a bezpečnost. Nezbytnost: kamera míří jen na vstup, nikoli na odpočinkovou místnost, záznam se drží krátce a přístup má jen omezený počet osob. Vyvážení: zaměstnanci a návštěvníci jsou jasně informováni, nejde o skryté sledování, záznam se používá jen při incidentu. Oprávněný zájem pravděpodobně obstojí. Kdyby kamera snímala celý pracovní prostor se zvukem a záznam se používal k výkonovému hodnocení zaměstnanců, vyvážení by nejspíš neobstálo a mohlo by být potřeba DPIA.}

\priklad{Bezpečnostní logy e-shopu: oprávněný zájem na ochraně účtů a detekci útoků je reálný. Nezbytnost znamená logovat údaje, které opravdu pomáhají detekci, ne sbírat všechno navždy. Vyvážení znamená krátké retenční lhůty, omezený přístup, informování uživatelů a nepoužívat bezpečnostní logy potichu k marketingovému profilování.}

\priklad{Aktuální příklad z praxe - údaje z realit a investiční nabídky: ÚOOÚ v kontrolách ukončených v roce 2025 popsal situaci, kdy byly údaje klientů realitního zprostředkovatele staženy z informačního systému a použity pro investiční nebo finanční nabídky bez odpovídajícího právního titulu a transparentního informování. Tento případ je dobrý na oprávněný zájem: obchodní zájem oslovit potenciální klienty může existovat, ale neznamená automaticky oprávněné převzetí databáze získané pro jiný účel. Ke zkoušce: balancing test zkoumá i původ dat, rozumné očekávání člověka, změnu účelu a možnost námitky.}

\subsection*{Riziko, DPIA a proporcionalita}
GDPR je rizikově orientované. Čím vyšší riziko pro práva a svobody osob, tím silnější musí být opatření. Když zpracování pravděpodobně povede k vysokému riziku, provádí se \textbf{DPIA} - posouzení vlivu na ochranu osobních údajů podle čl. 35 GDPR. DPIA se dělá \textbf{před zahájením zpracování}, ne až po incidentu nebo po spuštění systému.

Typicky je DPIA potřeba u:
\begin{itemize}
  \item systematického a rozsáhlého profilování s právními nebo obdobně závažnými účinky,
  \item rozsáhlého zpracování zvláštních kategorií údajů,
  \item rozsáhlého systematického monitorování veřejně přístupných prostor,
  \item biometrické identifikace, rozsáhlého kamerového sledování, nových technologií, AI systémů nebo kombinování datových sad,
  \item zpracování, které může vést k diskriminaci, ztrátě kontroly nad daty, finanční újmě, poškození pověsti nebo omezení práv.
\end{itemize}

Obsah DPIA v praxi:
\begin{enumerate}
  \item popsat účel a zamýšlené zpracování,
  \item určit kategorie údajů, subjektů, příjemců, zdroje a dobu uložení,
  \item doložit právní titul a případně výjimku podle čl. 9,
  \item posoudit nezbytnost a přiměřenost,
  \item identifikovat rizika pro osoby,
  \item navrhnout opatření: minimalizace, pseudonymizace, šifrování, řízení přístupu, logování, omezení retence, lidský dohled,
  \item vyhodnotit zbytkové riziko.
\end{enumerate}
Pokud i po opatřeních zůstává vysoké riziko, musí správce před zahájením zpracování konzultovat ÚOOÚ.

Test proporcionality se dá vyložit jednoduše:
\begin{enumerate}
  \item \textbf{legitimní cíl} - sleduji oprávněný a jasný účel?
  \item \textbf{vhodnost} - pomůže zpracování cíle dosáhnout?
  \item \textbf{nezbytnost} - nejde to méně invazivně?
  \item \textbf{přiměřenost v užším smyslu} - nepřevažuje zásah do práv nad přínosem?
\end{enumerate}

\priklad{Kamerový systém ve škole: účel může být bezpečnost. Vhodnost ano, kamery mohou pomoci. Nezbytnost se ptá, zda nestačí méně kamer, kratší uchování nebo jen vstupy. Přiměřenost řeší, zda se nesledují šatny, toalety nebo žáci nepřiměřeně dlouho.}

\subsection*{Pověřenec, privacy by design a profilování}
\textbf{Pověřenec pro ochranu osobních údajů} (DPO) je osoba, která radí správci nebo zpracovateli, sleduje soulad s GDPR, školí, kontroluje, spolupracuje s ÚOOÚ a je kontaktním místem pro subjekty údajů. Není to "ten, kdo za všechno odpovídá" místo správce; odpovědnost za soulad má pořád správce nebo zpracovatel.

DPO je povinný zejména:
\begin{itemize}
  \item u orgánů veřejné moci a veřejných subjektů,
  \item když hlavní činnost spočívá v rozsáhlém pravidelném a systematickém monitorování osob,
  \item když hlavní činnost spočívá v rozsáhlém zpracování zvláštních kategorií údajů nebo údajů podle čl. 10.
\end{itemize}
DPO musí být odborně způsobilý, nezávislý ve výkonu své funkce, nesmí být za plnění úkolů sankcionován a nemá se dostat do střetu zájmů. Malý e-shop DPO obvykle nepotřebuje, nemocnice nebo obec ano.

\textbf{Privacy by design a by default} znamená, že ochrana údajů se zabuduje už do návrhu systému a výchozí nastavení je co nejšetrnější: sbírat méně dat, kratší retence, omezené role, auditní logy, pseudonymizace, oddělení prostředí, výchozí neveřejnost profilu.

\textbf{Profilování} je automatizované zpracování osobních údajů k hodnocení osobních aspektů člověka, například ekonomické situace, zdraví, preferencí, chování, spolehlivosti nebo pohybu. Profilování není vždy zakázané, ale u rozsáhlého a systematického profilování s významnými dopady typicky roste riziko, může být potřeba DPIA a u čistě automatizovaného rozhodování s právními nebo obdobně závažnými účinky se uplatní čl. 22 GDPR.

\priklad{Nemocnice chce nasadit AI systém pro podporu diagnostiky. Zpracovává zdravotní údaje, tedy zvláštní kategorii. Potřebuje titul podle čl. 6 a výjimku podle čl. 9, DPO je prakticky povinný kvůli rozsáhlému zpracování citlivých údajů, před spuštěním se dělá DPIA a opatření jako pseudonymizace tréninkových dat, omezení přístupů, logování a lidský dohled nad rozhodnutím.}

\priklad{Aktuální příklad z praxe - AI modely a GDPR: EDPB ve stanovisku 28/2024 řeší, kdy lze model považovat za anonymní, kdy může být oprávněný zájem právním titulem pro vývoj nebo nasazení AI a jaké následky může mít nezákonné použití osobních údajů při trénování. Nejde o zákaz AI, ale o požadavek doložit účel, minimalizaci, právní titul, rizika, práva subjektů a vliv případného nezákonného zdroje dat. Ke zkoušce: u AI se neptám jen "je model užitečný", ale i "jaká data ho vytvořila, zda jsou osoby identifikovatelné a jaké dopady má výstup na člověka".}

\subsection*{Porušení zabezpečení a práva subjektů}
Porušení zabezpečení osobních údajů je únik, ztráta, zničení, změna nebo neoprávněný přístup. Správce oznamuje dozorovému úřadu zpravidla do 72 hodin od okamžiku, kdy se o porušení dozvěděl, pokud je pravděpodobné riziko pro práva a svobody. Subjektům údajů oznamuje bez zbytečného odkladu až tehdy, pokud je vysoké riziko. Pokud je riziko nepravděpodobné, nemusí se hlásit ÚOOÚ, ale porušení se má interně zdokumentovat.

Při hodnocení rizika se řeší povaha údajů, rozsah, identifikovatelnost, zranitelnost osob, zabezpečení, možnost zneužití a následky. Únik hashovaných hesel bude jiný, pokud jsou hesla silně solená a hashovaná moderním algoritmem, než pokud jde o plaintext hesla se jmény a e-maily. Pokud je organizace zároveň regulovaná podle ZKB, může vedle GDPR běžet i samostatné hlášení kybernetického incidentu.

Subjekty údajů mají práva podle čl. 12-22 GDPR: informace, přístup, opravu, výmaz, omezení zpracování, přenositelnost, námitku a nebýt předmětem čistě automatizovaného rozhodování s právními nebo obdobně významnými účinky. Žádost se vyřizuje zpravidla do jednoho měsíce, s možností prodloužení o dva měsíce u složitých případů.

\subsubsection*{Námitka vs výmaz}
\textbf{Námitka} podle čl. 21 míří hlavně na zpracování založené na oprávněném zájmu nebo veřejném zájmu. U přímého marketingu je zvlášť silná: pokud subjekt vznese námitku proti přímému marketingu, údaje už pro tento účel nesmí být zpracovávány. U jiného oprávněného zájmu může správce pokračovat jen tehdy, pokud prokáže závažné oprávněné důvody, které převažují nad zájmy a právy subjektu, nebo pokud data potřebuje pro právní nároky.

\textbf{Výmaz} podle čl. 17 znamená odstranění údajů, ale není absolutní. Neuplatní se, pokud správce údaje potřebuje například ke splnění právní povinnosti, pro výkon nebo obhajobu právních nároků, veřejný zájem, archivaci nebo svobodu projevu. Prakticky: zákazník může požádat o výmaz marketingového profilu, ale faktury se nevymažou, pokud je správce musí držet podle účetních a daňových předpisů.

\subsection*{ÚOOÚ}
ÚOOÚ je nezávislý dozorový úřad. Kontroluje dodržování GDPR a českého zákona, řeší stížnosti, ukládá opatření a pokuty, schvaluje některé nástroje, konzultuje DPIA, vydává metodiku a spolupracuje v Evropském sboru pro ochranu osobních údajů (EDPB). Nezávislost je klíčová, protože úřad dohlíží i na veřejnou moc.

Sankce mohou dosáhnout 20 mil. eur nebo 4 procent celosvětového ročního obratu u závažnějších porušení, případně 10 mil. eur nebo 2 procenta u méně závažných porušení. V přeshraničních věcech funguje mechanismus one-stop-shop: hlavní dozorový úřad je typicky úřad státu hlavní provozovny správce, ostatní dotčené úřady spolupracují a EDPB pomáhá sjednocovat výklad.

\subsection*{Předávání mimo EHP}
Předání osobních údajů mimo Evropský hospodářský prostor je možné jen při zajištění srovnatelné ochrany. Nejjednodušší je rozhodnutí Komise o odpovídající ochraně. Pokud není, používají se standardní smluvní doložky, závazná podniková pravidla nebo výjimky pro specifické situace. Po Schrems II je potřeba u standardních doložek posoudit i právo cílové země a praktická rizika přístupu veřejné moci, tedy transfer impact assessment.

Pro USA existuje EU-US Data Privacy Framework jako rozhodnutí o odpovídající ochraně pro certifikované subjekty. K 18. 6. 2026 platí, ale po rozhodnutí Tribunálu ve věci Latombe z 3. 9. 2025 běží odvolací rovina u Soudního dvora, takže u zkoušky je dobré říct: rámec se používá, ale politicko-právní nejistota kolem transatlantických přenosů trvá.

\subsection*{Modelové situace}
\begin{enumerate}
  \item \textbf{E-shop}: doručení objednávky je plnění smlouvy, faktury právní povinnost, bezpečnostní logy oprávněný zájem, newsletter novému zájemci souhlas, nabídka obdobného zboží stávajícímu zákazníkovi může spadat pod soft opt-in. Námitka proti marketingu zastaví marketing, ale nevymaže faktury.
  \item \textbf{Kamerový systém v obci}: veřejná moc nebo oprávněný zájem podle konkrétního režimu, jasný účel, informační cedule, omezená doba uložení, záběr jen potřebných míst, případně DPIA při rozsáhlém systematickém monitorování.
  \item \textbf{HR screening uchazečů}: nelze bez dalšího sbírat vše z internetu. Účel musí souviset s pozicí, rozsah musí být přiměřený, uchazeč má být informován a zvláštní kategorie údajů jsou vysoce rizikové.
  \item \textbf{AI profilování klientů banky}: profilování kvůli prevenci podvodů může mít oprávněný zájem nebo právní povinnost, ale je potřeba posoudit rizika, vysvětlitelnost, lidský dohled, minimalizaci a často DPIA.
\end{enumerate}

\priklad{Avast/Jumpshot: ÚOOÚ uložil Avast Software pokutu 351 mil. Kč za neoprávněné zpracování údajů o prohlížení uživatelů. Klíčová učební pointa je pseudo vs anonym: odstranění přímých identifikátorů a nahrazení identifikátorem nemusí být anonymizace, pokud lze osoby rozumně znovu identifikovat nebo profilovat. Městský soud v Praze v roce 2025 potvrdil porušení, ale zrušil část rozhodnutí ohledně výše pokuty kvůli odůvodnění; věc tím dobře ukazuje, že "anonymizováno" nestačí jen tvrdit.}

\zkouska{U GDPR odpovědi nikdy nezačínej "potřebujeme souhlas". Začni účelem, pak právním titulem, zásadami, minimalizací, rizikem a dokumentací. U oprávněného zájmu řekni tři kroky LIA; u rizikového zpracování řekni DPIA před spuštěním; u citlivých údajů řekni čl. 6 plus čl. 9.}

\zdroje{GDPR 2016/679; zákon o zpracování osobních údajů (č. 110/2019 Sb.); zákon o některých službách informační společnosti (č. 480/2004 Sb.); metodiky ÚOOÚ k DPIA a DPO; EDPB Guidelines 1/2024 k oprávněnému zájmu; věc Breyer C-582/14; Schrems II C-311/18; věc Latombe T-553/23 a navazující odvolání; ÚOOÚ a Městský soud v Praze ve věci Avast/Jumpshot; přiložené PDF a vložený text.}

\newpage
\fi
